%% 02-solution.tex — The algorithmic solution

\section{The Solution}

\noindent\textbf{Replace the politician with an algorithm.}

\medskip
\noindent
The algorithm works like splitting a pizza: make one cut through the
middle of the state, dividing it into two equal-population halves.
Cut each half in half.
Repeat until the state has as many pieces as it has congressional seats.
Every cut is guided by one rule---make it as short as possible---so
districts follow natural geographic features instead of zigzagging to
capture partisan voters.

\medskip
\noindent
The algorithm uses U.S.\ Census population counts and tract boundary
geometries, both of which are public data~\cite{census2020pl}.
It never sees voter registration files, election results, or incumbents'
home addresses.
\textbf{It cannot gerrymander because it lacks the information
required to do so.}
This is not a policy choice that a future legislature could reverse;
it is a structural property of the algorithm.

\medskip
\noindent
The algorithm, called \textit{ApportionRegions}, is a direct extension of
\textit{Huntington-Hill}---the mathematical method Congress already uses to
apportion seats among states, mandated by statute since
1941~\cite{huntingtonHill1941}.
ApportionRegions extends the same principle from apportionment to intrastate
district drawing; the two are related in approach but distinct in statutory
scope.
Extending the principle to boundary design requires one sentence of statute.

\medskip
\noindent
Two independent verifiers running the algorithm on the same Census data
produce byte-identical district assignments, auditable through a
SHA-256 hash chain.
Any citizen can verify the result in 15--20 minutes end-to-end,
including data download.
